%! Sinusoidal Functions — Unit 6, Lesson 6.6 — Student Packet Cover Sheet
\documentclass[10pt]{article}
\usepackage{precalculus-article}
\usepackage{precalculus-boxes}
\usepackage{ltablex}
\keepXColumns

\begin{document}

% Full-bleed navy banner
\begin{tikzpicture}[remember picture, overlay]
  \fill[navy] (current page.north west)
    rectangle ([yshift=-0.9in]current page.north east);
\end{tikzpicture}
\vspace{-0.8in}

\begin{center}
  {\color{white}\textbf{\LARGE Precalculus}}\\[4pt]
  {\color{skymid}\large\textbf{Unit 6: Trigonometric Functions}}\\[6pt]
  {\color{white}\textbf{\large Lesson 6.6 \quad Sinusoidal Functions}}
\end{center}

\vspace{0.22in}
\namedateperiod
\vspace{0.18in}

\begin{learningtargetbox}
\small
\begin{enumerate}[leftmargin=1.5em, itemsep=5pt]
  \item I can define a sinusoidal function and explain the relationship
        $\cos\theta = \sin(\theta + \pi/2)$.
  \item I can determine the period, frequency, amplitude, and midline of a
        sinusoidal function from its graph or a verbal description.
  \item I can describe the symmetry of $y = \sin\theta$ (odd) and $y = \cos\theta$
        (even) and explain what that means graphically.
\end{enumerate}
\end{learningtargetbox}

\vspace{0.14in}

\begin{tocbox}
\small
\renewcommand{\arraystretch}{1.5}
\begin{tabularx}{\linewidth}{@{} c l X @{}}
\rowcolor{sky}
\textbf{\#} & \textbf{Component} & \textbf{Description} \\
\midrule
1 & Warm-Up        & Review of sine and cosine graphs from Lesson 3.4 \\
2 & Guided Notes   & Characteristics of sinusoidal functions \\
3 & Group Activity & Identifying characteristics at three tiers \\
4 & Exit Ticket    & Amplitude, midline, and frequency from description \\
5 & Homework       & Mixed practice and AP-style problem \\
\end{tabularx}
\end{tocbox}

\end{document}
